\chapter*{Preface}
% BEGIN CURATED SUBJECT INDEX
\index{unbounded operator}
\index{Hilbert space}
\index{spectral measure}
\index{spectrum!point}
\index{connection coefficient}
\index{Breit--Wigner approximation}
\index{exact WKB}
\index{Zel'dovich regularization}
\index{rigged Hilbert space}
% END CURATED SUBJECT INDEX
\addcontentsline{toc}{chapter}{Preface}
This textbook grew out of our studies of quantum resonances at the intersection of resurgence and non-self-adjoint spectral theory. These investigations gradually led us to the scattering-theoretic viewpoint that organizes this book. Our aim is not merely to present this viewpoint, but also to develop the mathematical structures on which it rests and to examine the physical phenomena through which it becomes concrete. The book is therefore intended both as a self-contained introduction to the relevant mathematics and physics and as a systematic account of the scattering-theoretic perspective that emerged from our own research.


Resonances are often taught in reverse order.  One first sees a
Breit--Wigner denominator or a complex eigenvalue, learns several devices for
making its wave function finite, and only later discovers that each device
continues a scattering problem across a cut.  That route is efficient when
one number is required.  It is much less satisfactory when the questions are
why the number is invariant, what space contains the associated state, how
the continuum participates, or what changes at a threshold.

This book follows the opposite route.  We begin with standard quantum
mechanics, including the pieces that are usually suppressed by notation:
operator domains, limiting procedures, spectral measures, and comparison
dynamics.  We then build scattering data and analytically continue the
global connection coefficient.  A resonance is defined there.  Complex
scaling, Zel'dovich regularization, rigged Hilbert spaces, and effective
Hamiltonians enter later because they solve particular representation or
computational problems.  Exact WKB is the thread that keeps these changes of
representation tied to the same analytic object.

The result is a monograph rather than a survey.  A citation is never meant to
replace an equation that the reader needs in order to continue a calculation.
Historical sources and alternative hypotheses are recorded in the
bibliography, but the main line supplies the definitions, derivations, and
checks.  Results from our papers have consequently been disassembled and
rebuilt in prerequisite order.  Introductory material that occurred in more
than one paper appears once in the canonical chapter; a later Laboratory
points back to it and starts at the first genuinely new step.


%Parts of the exposition, detailed derivations, and figures in this book are based on, and substantially extend, material first presented in our research articles. They have been reorganized and expanded here to form a self-contained pedagogical development. References to the original publications are given in the bibliographical notes and bibliography.

\section*{How to read the Stages}

Each Stage begins with a reader contract and ends with problems.  The contract
states an operation the reader should be able to perform, not a list of topics
to have seen.  A first reading may follow Stages I--IV linearly and then choose
among the nuclear, exceptional-point, and black-hole branches.  A calculation
oriented reading can instead begin at a Laboratory and use the margin of
cross-references to retrieve its prerequisites.

The logical dependency is
\begin{gather*}
  \boxed{\text{domain and spectrum}}
  \quad\Downarrow\quad
  \boxed{\text{scattering boundary data}}
  \\
  \boxed{\text{global connection datum}}
  \quad\Downarrow\quad
  \boxed{\text{realizations and applications}}.
\end{gather*}
The arrows are one-way in the derivation but two-way in practice.  A numerical
complex-scaled spectrum, for example, can reveal a candidate pole; the
connection formulation tells us which variations must leave it fixed before
it is accepted as physical.

\section*{What a Laboratory contains}

A Laboratory is not a paper reprint.  It has four components:
\begin{enumerate}[label=\arabic*.]
  \item the input Hamiltonian or master equation and a declared domain;
  \item a derivation with enough intermediate algebra to reproduce it;
  \item a stability or conservation check that can falsify the result; and
  \item a checkpoint stating what has actually been learned.
\end{enumerate}
Some Laboratories retain figures and detailed calculations developed in the
source papers.  Their role here is different: they are worked examples of a
method already constructed, so paper-specific motivation and repeated mini-
reviews have been removed.

\section*{Prerequisites and conventions}

The assumed physics is an undergraduate course in quantum mechanics and
familiarity with ordinary differential equations and complex analysis.
Functional analysis, scattering theory, resurgence, general relativity, and
few-body reaction theory are developed when first needed.  A reader already
fluent in these subjects can use the opening problems of each Stage as a
diagnostic and skip proofs only after checking the stated hypotheses.

We use the time dependence $\rme^{-\rmi Et/\hbar}$, so a decaying pole has
$\im E<0$.  A differential expression and the operator defined by it are
distinguished whenever boundary conditions or domains matter.  The symbol
$H_\theta$ denotes an analytically deformed operator, not an arbitrary
non-Hermitian approximation.  These choices are collected once in the
conventions chapter and are not redefined in later Laboratories.

\section*{Acknowledgment of provenance}
% BEGIN CURATED SUBJECT INDEX
\index{CDCC}
\index{lithium nuclei}
\index{beryllium nuclei}
% END CURATED SUBJECT INDEX

The exact-WKB, complex-scaling, continuum, radial, exceptional-point, and
black-hole calculations build on the authors' research program and retain
its detailed algebra where that is pedagogically useful.  The few-body
chapters connect the same spectral viewpoint to CDCC and to calculations of
weakly bound lithium and beryllium-core systems.  Bibliographic citations
identify original results; within the book, however, their content is used as
raw material for one continuous derivation.


%This work was partially supported by Japan Society for the Promotion of Science (JSPS) Grant-in-Aid for Scientific Research Grant Numbers JP25K17402 (O.M.), and 25H01267 (S.O.).
%This work was also supported in part by the COREnet project of RCNP, The University of Osaka.
%O.M.\ acknowledges the RIKEN Special Postdoctoral Researcher Program and RIKEN FY2025 Incentive Research Projects.
